\documentclass[10pt,a4paper]{amsart}
\usepackage{amsmath,amssymb,amsthm,graphicx}
\usepackage{verbatim}
\newenvironment{reference}{\comment}{\endcomment}
%\newenvironment{reference}{}{}
\newenvironment{slogan}{\comment}{\endcomment}
\newenvironment{history}{\comment}{\endcomment}

% For commutative diagrams we use Xy-pic
\usepackage[all]{xy}

% We use 2cell for 2-commutative diagrams.
\xyoption{2cell}
% Theorem environments.
%
\theoremstyle{plain}
\newtheorem{theorem}[subsection]{Theorem}
\newtheorem{proposition}[subsection]{Proposition}
\newtheorem{lemma}[subsection]{Lemma}

\theoremstyle{definition}
\newtheorem{definition}[subsection]{Definition}
\newtheorem{example}[subsection]{Example}
\newtheorem{exercise}[subsection]{Exercise}
\newtheorem{situation}[subsection]{Situation}

\theoremstyle{remark}
\newtheorem{remark}[subsection]{Remark}
\newtheorem{remarks}[subsection]{Remarks}

\numberwithin{equation}{subsection}

% Macros
%
\def\lim{\mathop{\mathrm{lim}}\nolimits}
\def\colim{\mathop{\mathrm{colim}}\nolimits}
\def\Spec{\mathop{\mathrm{Spec}}}
\def\Hom{\mathop{\mathrm{Hom}}\nolimits}
\def\Ext{\mathop{\mathrm{Ext}}\nolimits}
\def\SheafHom{\mathop{\mathcal{H}\!\mathit{om}}\nolimits}
\def\SheafExt{\mathop{\mathcal{E}\!\mathit{xt}}\nolimits}
\def\Sch{\mathit{Sch}}
\def\Mor{\mathop{\mathrm{Mor}}\nolimits}
\def\Ob{\mathop{\mathrm{Ob}}\nolimits}
\def\Sh{\mathop{\mathit{Sh}}\nolimits}
\def\NL{\mathop{N\!L}\nolimits}
\def\CH{\mathop{\mathrm{CH}}\nolimits}
\def\proetale{{pro\text{-}\acute{e}tale}}
\def\etale{{\acute{e}tale}}
\def\QCoh{\mathit{QCoh}}
\def\Ker{\mathop{\mathrm{Ker}}}
\def\Im{\mathop{\mathrm{Im}}}
\def\Coker{\mathop{\mathrm{Coker}}}
\def\Coim{\mathop{\mathrm{Coim}}}

% Boxtimes
%
\DeclareMathSymbol{\boxtimes}{\mathbin}{AMSa}{"02}

%
% Macros for moduli stacks/spaces
%
\def\QCohstack{\mathcal{QC}\!\mathit{oh}}
\def\Cohstack{\mathcal{C}\!\mathit{oh}}
\def\Spacesstack{\mathcal{S}\!\mathit{paces}}
\def\Quotfunctor{\mathrm{Quot}}
\def\Hilbfunctor{\mathrm{Hilb}}
\def\Curvesstack{\mathcal{C}\!\mathit{urves}}
\def\Polarizedstack{\mathcal{P}\!\mathit{olarized}}
\def\Complexesstack{\mathcal{C}\!\mathit{omplexes}}
% \Pic is the operator that assigns to X its picard group, usage \Pic(X)
% \Picardstack_{X/B} denotes the Picard stack of X over B
% \Picardfunctor_{X/B} denotes the Picard functor of X over B
\def\Pic{\mathop{\mathrm{Pic}}\nolimits}
\def\Picardstack{\mathcal{P}\!\mathit{ic}}
\def\Picardfunctor{\mathrm{Pic}}
\def\Deformationcategory{\mathcal{D}\!\mathit{ef}}
\usepackage{cleveref}
\crefname{originalequation}{original equation}{original equations}
\Crefname{originalequation}{Original equation}{Original equations}
\makeatletter
\newlabel{equation-stratify}{{096X}{source}{}{sourcepointer-1}{}}
\newlabel{equation-stratify@cref}{{[originalequation][0][]096X}{[1][0][]source}{}{}{}}
\makeatother
\title{Every finite constructible stratification of an affine scheme is refined by the stratification attached to a finite set of elements}
\author{The Stacks Project Authors}
\date{stacks-project:proetale.tex:local-snapshot-20261012; GNU FDL 1.2 or later}
\begin{document}
\maketitle
\noindent\textit{Editorial note.} All delivered mathematical text is the author's own, copied from the pinned Stacks Project chapter file proetale.tex; no statement or proof step was rewritten, reordered or removed.\par\medskip
\noindent\textit{Editorial note.} The counted proof is the author's complete argument: the presentation of the constructible strata through quasi-compact opens and their complements, the replacement of those opens by finitely many basic opens, and the verification that the vanishing-pattern stratification refines the given one.\par\medskip
\noindent\textit{Editorial note.} The label lemma-refine names the counted statement, and the proof environment immediately following it in the source is delivered verbatim as its proof.\par\medskip
\noindent\textit{Editorial note.} The counted statement is carried by the author's own lemma environment in the chapter source.\par\medskip
\noindent\textit{Editorial note.} Pre-open gate: python operations/agent-authoring/tools/precard\_gate.py --source proetale.tex --label lemma-refine --start 770 --end 787 --local-id 28006 returned RESULT=OK (L28006 free, occupancy FREE, public interval clear). Candidate taken from sources/stacks-proetale/candidates-halfB.json.\par\medskip
\noindent\textit{Editorial note.} The statement and proof are copied from the pinned chapter file; the cited inputs are kept as external pointers with their canonical Stacks tags.\par\medskip
\noindent\textit{Allowed standard prerequisites (editorial).}\begin{itemize}
\item Presentation of a constructible subset as a finite union of differences of quasi-compact opens. Topology, the description of constructible subsets of a spectral space as finite unions of locally closed subsets built from quasi-compact opens. Supplies the writing of each stratum T\_i as a finite union of intersections U\_\{i,j\} cap (V\_\{i,j\})\textasciicircum{}c.
\item Description of quasi-compact opens by basic opens. Algebra and Schemes, the description of a quasi-compact open of the spectrum of a ring as a finite union of basic opens D(f). Supplies the elements f\_\{i,j,k\} and g\_\{i,j,l\} whose basic opens cut out the U\_\{i,j\} and V\_\{i,j\}.
\item The stratification attached to a finite set of elements. Same chapter, the definition of the stratification whose strata are labelled by the vanishing pattern of the elements of a finite set. Supplies the stratification that the construction refines the given one.
\end{itemize}\medskip
\ifdefined\hypertarget\hypertarget{sourcepointer-1}{}\fi
\noindent\textit{Source reference.} Original equation 096X: sources/stacks-proetale/pinned-original-source.tex, source lines 765--765. Prerequisite pointer invoked at proof line(s) 774; it names a standard fact of the same chapter that the counted argument uses, and every step of the counted proof built on it is retained verbatim in the delivered body.\par\medskip
\begin{lemma}
\label{lemma-refine}
Let $X = \Spec(A)$ as above. Given any finite stratification
$X = \coprod T_i$ by constructible subsets, there exists a finite
subset $E \subset A$ such that the stratification (\ref{equation-stratify})
refines $X = \coprod T_i$.
\end{lemma}

\begin{proof}
We may write $T_i = \bigcup_j U_{i, j} \cap V_{i, j}^c$ as a finite union
for some $U_{i, j}$ and $V_{i, j}$ quasi-compact open in $X$.
Then we may write $U_{i, j} = \bigcup D(f_{i, j, k})$ and
$V_{i, j} = \bigcup D(g_{i, j, l})$. Then we set
$E = \{f_{i, j, k}\} \cup \{g_{i, j, l}\}$. This does the job, because
the stratification (\ref{equation-stratify}) is the one whose strata are
labeled by the vanishing pattern of the elements of $E$ which
clearly refines the given stratification.
\end{proof}

\end{document}
